\documentclass[10pt]{article}
\usepackage[utf8]{inputenc}
\usepackage[T5]{fontenc} % Sử dụng encoding T5 để hỗ trợ tiếng Việt
\usepackage[vietnam]{babel} % Gói hỗ trợ tiếng Việt
\usepackage{amsmath}
\usepackage{amsfonts}
\usepackage{amssymb}
\usepackage[version=4]{mhchem}
\usepackage{stmaryrd}
\usepackage{bbold}

\begin{document}
ký hiệu trống. Vì lý do thuận tiện sau này, chúng tôi sử dụng " $b$ " để chỉ ký hiệu trống, thay vì $\#$. Theo cách tương tự, giả sử rằng kết quả của một hàm $\square_{1}^{\mathrm{QP}}$ $f$ bao gồm một phần quan trọng, có ý nghĩa và phần "rác" còn lại, là tàn dư của quá trình tính toán. Vì chúng tôi chỉ sử dụng các qubit ( $|0\rangle$ và $|1\rangle$ ) để biểu diễn đầu vào và đầu ra, làm thế nào chúng tôi có thể phân biệt giữa phần có ý nghĩa và phần rác? Để mô phỏng các QTMs bằng các hàm $\square_{1}^{\text {QP }}$, do đó chúng tôi cần "mã hóa" tất cả các ký hiệu băng thành các qubit.

Bài viết này giới thiệu sơ đồ mã hóa đơn giản sau đây. Chúng tôi đặt $\hat{0}=00, \hat{1}=01$, và $\hat{b}=10$. Chúng tôi cũng đặt $\hat{2}=11$ và $\hat{3}=10$ để sử dụng sau này. Bảng chữ cái đầu vào $\Sigma=\{0,1\}$ do đó được dịch sang $\{\hat{0}, \hat{1}\}$ và bảng chữ cái băng $\Gamma=\Sigma \cup\{b\}$ được mã hóa thành $\{\hat{0}, \hat{1}, \hat{b}\}$. Cho mỗi chuỗi nhị phân $s=s_{1} s_{2} \cdots s_{n}$ với $s_{i} \in\{0,1\}$ cho mọi chỉ số $i \in[n]$, một mã $\tilde{s}$ của $s$ biểu thị chuỗi $\hat{s}_{1} \hat{s}_{2} \cdots \hat{s}_{n} \hat{2}$, trong đó mục cuối cùng $\hat{2}$ phục vụ như một dấu hiệu kết thúc, đánh dấu cuối mã. Khi đó ta có $|\tilde{s}|=2|s|+2$. Là những ví dụ nhanh, ta có được $\widetilde{0110}=\hat{0} \hat{1} \hat{1} \hat{0} \hat{2}=0001010011$ và $\widetilde{01} \cdot \widetilde{11}=\hat{0} \hat{1} \hat{2} \hat{1} \hat{2}=000111010111$. Sau này, chúng tôi sẽ chỉ ra trong bổ đề 5.1 rằng việc mã hóa các chuỗi và giải mã các chuỗi đã mã hóa có thể được thực hiện bởi các hàm $\widehat{\square_{1}^{\mathrm{QP}}}$ thích hợp.

Một khó khăn khác đến từ việc các hàm $\widehat{\square_{1}^{\mathrm{QP}}}$ không thể mở rộng các qubit của chúng. Lưu ý rằng một QTM được thiết kế để tự do sử dụng không gian lưu trữ bổ sung bằng cách di chuyển đầu băng của nó đơn giản đến các ô băng mới, trống rỗng vượt ra ngoài khu vực đầu vào, nơi mà một đầu vào ban đầu được viết. Để mô phỏng một QTM như vậy, chúng ta cần mô phỏng toàn bộ các hoạt động của nó được thực hiện trong không gian băng mở rộng tự do. Ngược lại, mọi hàm $\widehat{\square_{1}^{\mathrm{QP}}}$ đều bảo toàn chiều theo Bổ đề 3.4(4), do đó số lượng qubit đầu vào phải khớp với số lượng qubit đầu ra. Nếu chúng ta muốn mô phỏng không gian lưu trữ bổ sung của QTM, thì chúng ta cần cấp cho hàm $\widehat{\square_{1}^{\mathrm{QP}}}$ mục tiêu cùng lượng qubit bổ sung để sử dụng ngay từ đầu.

Chúng tôi giải quyết vấn đề thứ hai này bằng cách mở rộng mỗi đầu vào của các hàm lượng tử bằng cách thêm các số 0 bổ sung có độ dài liên quan đến thời gian chạy của QTM. Đối với bất kỳ đa thức $p$ nào và bất kỳ hàm lượng tử $g$ nào trên $\mathcal{H}_{\infty}$, chúng tôi định nghĩa $\left|\phi^{p}(x)\right\rangle=\left|0^{|x|} 1\right\rangle\left|0^{p(|x|)} 10^{11 p(|x|)+6} 1\right\rangle|x\rangle$ và $\left|\phi_{g}^{p}(x)\right\rangle=g\left(\left|\phi^{p}(x)\right\rangle\right)$ cho mọi chuỗi $x \in\{0,1\}^{*}$. Tương tự, đối với bất kỳ hàm $f$ nào trên $\{0,1\}^{*}$, chúng tôi đặt $\left.\left|\phi^{p, f}(x)\right\rangle=\left|\widetilde{\left.0^{|f(x)|}\right\rangle}\right| 0^{|f(x)|+1} 1\right\rangle\left|\phi^{p}(x)\right\rangle$ và $\left|\phi_{g}^{p, f}(x)\right\rangle=g\left(\left|\phi^{p, f}(x)\right\rangle\right)$.

Bất kỳ hàm FBQP nào cũng nhận các chuỗi đầu vào cổ điển và tạo ra các chuỗi cổ điển là kết quả của một máy Turing lượng tử chạy trong thời gian đa thức với xác suất cao. Trái lại, các hàm $\square_{1}^{\mathrm{QP}}$ của chúng tôi $f$ là để biến đổi mỗi trạng thái lượng tử trong $\mathcal{H}_{\infty}$ thành một trạng thái khác trong $\mathcal{H}_{\infty}$. Để có được các chuỗi đầu ra cổ điển, chúng ta cần quan sát kết quả của $f$.

Định lý chính (Định lý 4.1) khẳng định đại khái như sau: đối với bất kỳ hàm FBQP $f$ nào, luôn tồn tại một hàm lượng tử $g$ trong $\widehat{\square_{1}^{\mathrm{QP}}}$ sao cho, khi chúng ta quan sát $|f(x)|$ qubit đầu tiên của kết quả $g\left(\left|\phi^{p}(x)\right\rangle\right)(= \left|\phi_{g}^{p}(x)\right\rangle$ ) của $g$ trên đầu vào $\left|\phi^{p}(x)\right\rangle$, chúng ta nhận được $f(x)$ một cách chính xác với xác suất cao. Lưu ý rằng $g\left(\left|\phi^{p}(x)\right\rangle\right)$ cũng chứa các qubit bổ sung, gọi là các qubit "rác", mà không được quan sát trong quá trình tính toán $f(x)$. Những qubit rác này thực sự là tàn dư của quá trình tính toán $f$. Tuy nhiên, có thể loại bỏ những qubit đó bằng cách một phần đảo ngược toàn bộ quá trình tính toán, nếu chúng ta biết kích thước đầu ra $|f(x)|$ trước khi tính toán (bằng cách mở rộng $\left|\phi^{p}(x)\right\rangle$ thành $\left|\phi^{p, f}(x)\right\rangle$ và $\left|\phi_{g}^{p}(x)\right\rangle$ thành $\left|\phi_{g}^{p, f}(x)\right\rangle$ ).

Định lý 4.1 (Định lý chính) Cho $f$ là một hàm có giới hạn đa thức trên $\{0,1\}^{*}$. Ba phát biểu sau đây là tương đương về mặt logic.

\begin{enumerate}
  \item Hàm $f$ thuộc FBQP .
  \item Đối với bất kỳ hằng số $\varepsilon \in[0,1 / 2)$ nào, tồn tại một hàm lượng tử $g$ trong $\widehat{\square_{1}^{\mathrm{QP}}}$ và một đa thức $p$ sao cho $|f(x)| \leq p(|x|)$ và $\left\|\left\langle f(x) \mid \phi_{g}^{p}(x)\right\rangle\right\|^{2} \geq 1-\varepsilon$ cho tất cả $x \in\{0,1\}^{*}$.
  \item Đối với bất kỳ hằng số $\varepsilon \in[0,1 / 2)$ nào, tồn tại một hàm lượng tử $g$ trong $\widehat{\square_{1}^{\mathrm{QP}}}$ và một đa thức $p$ sao cho $|f(x)| \leq p(|x|)$ và $\left|\left\langle\Psi_{f(x)} \mid \phi_{g}^{p, f}(x)\right\rangle\right|^{2} \geq 1-\varepsilon$ cho tất cả $x \in\{0,1\}^{*}$, trong đó $\left|\Psi_{f(x)}\right\rangle= |f(x)\rangle\left|\left(0^{|f(x)|+1} 1\right)^{2}\right\rangle\left|\phi^{p}(x)\right\rangle$.
\end{enumerate}

Trong các phát biểu 2-3 của Định lý 4.1, $\left\langle f(x) \mid \phi_{g}^{p}(x)\right\rangle$ là một vector khác rỗng trong khi $\left\langle\Psi_{f(x)} \mid \phi_{g}^{p, f}(x)\right\rangle$ chỉ là một đại lượng vô hướng vì $\ell\left(|f(x)\rangle\left|\phi^{p}(x)\right\rangle\right)=\ell\left(\left|\phi_{g}^{p, f}(x)\right\rangle\right)$.

Ở đây, chúng tôi muốn chứng minh định lý chính, Định lý 4.1. Vì lý do chiến lược, chúng tôi chia định lý thành hai bổ đề kỹ thuật, Bổ đề 4.2 và 4.3 Để trình bày các bổ đề, chúng tôi cần giới thiệu thêm thuật ngữ cho QTMs. Việc thiết kế các QTMs nhiều băng dễ dàng hơn so với các QTM băng đơn. Tuy nhiên, vì các QTM nhiều băng có thể được mô phỏng bởi các QTM băng đơn bằng cách dịch nhiều băng thành nhiều dải của một băng đơn, nên để chứng minh Định lý 4.1, chỉ cần tập trung vào các QTM băng đơn.

Một QTM băng đơn được gọi là ở dạng chuẩn nếu $\delta\left(q_{f}, \sigma\right)=\left|q_{0}\right\rangle|\sigma\rangle|R\rangle$ đúng với bất kỳ ký hiệu băng nào\\
$\sigma \in \Gamma$. Nếu một QTM dừng lại trong một chồng chất các cấu hình cuối cùng mà đầu băng quay trở lại ô bắt đầu, thì ta gọi máy như vậy là tĩnh. Tham khảo [5] để biết các tính chất cơ bản của chúng. Vì tiện lợi, ta gọi một QTM là bảo thủ nếu nó được hình thành tốt, tĩnh và ở dạng chuẩn. Hơn nữa, một QTM được hình thành tốt được gọi là thuần túy nếu hàm chuyển đổi của nó thỏa mãn yêu cầu cụ thể sau: đối với mỗi cặp $(p, \sigma) \in Q \times \Gamma, \delta(p, \sigma)$ có dạng $\delta(p, \sigma)=e^{i \theta}|q, \tau, d\rangle$ hoặc $\delta(p, \sigma)=\cos \theta|q, \tau, d\rangle+\sin \theta\left|q^{\prime}, \tau^{\prime}, d^{\prime}\right\rangle$ cho một số $\theta \in[0,2 \pi)$ và hai bộ ( $q, \tau, d$ ) và ( $q^{\prime}, \tau^{\prime}, d^{\prime}$ ) khác biệt. Bernstein và Vazirani [5] khẳng định rằng bất kỳ QTM bảo thủ, thời gian đa thức, băng đơn nào $M$ cũng có thể được mô phỏng bởi một QTM bảo thủ, thời gian đa thức, băng đơn, thuần túy thích hợp $M^{\prime}$. Ngoài ra, khi $M$ có biên độ $\tilde{\mathbb{C}}$, thì $M^{\prime}$ cũng vậy.

Hãy nhớ rằng bất kỳ chuỗi đầu ra nào của một QTM đều bắt đầu từ ô bắt đầu và kéo dài về bên phải cho đến ký hiệu trống đầu tiên mặc dù có thể có các ký hiệu băng còn lại chưa được xóa ở các phần khác của băng. Vì tiện lợi, một QTM $M$ được gọi là có đầu ra sạch nếu, khi $M$ dừng lại, không có ký hiệu nào không trống xuất hiện ở vùng bên trái của chuỗi đầu ra, tức là vùng bao gồm tất cả các ô được đánh chỉ số bởi các số nguyên âm. Lấy một đa thức $p$ giới hạn thời gian chạy của $M$ trên mọi đầu vào. Vì $M$ dừng lại trong nhiều nhất $p(|x|)$ bước trên mọi trường hợp $x \in\{0,1\}^{*}$, nên chúng ta chỉ cần chú ý đến vùng băng thiết yếu của nó bao gồm tất cả các ô băng được đánh chỉ số từ $-p(|x|)$ đến $+p(|x|)$. Trong thực tế, chúng ta định nghĩa lại một cấu hình $\gamma$ của $M$ trên đầu vào $x$ có độ dài $n$ là một bộ ba ( $q, h, \sigma_{1} \cdots \sigma_{2 p(n)+1}$ ), trong đó $p \in Q, h \in \mathbb{Z}$ với $-p(n) \leq h \leq p(n)$, và $\sigma_{1}, \cdots, \sigma_{2 p(n)+1} \in\{0,1, b\}$ sao cho, với mọi chỉ số $i \in[2 p(n)+1], \sigma_{i}$ là một ký hiệu băng được viết tại ô được đánh chỉ số bởi $i-p(n)-1$. Lưu ý rằng ô bắt đầu nằm ở giữa $\sigma_{1} \cdots \sigma_{2 p(n)+1}$. Vì tiện lợi, ta tiếp tục sửa đổi khái niệm cấu hình bằng cách chia vùng băng thiết yếu thành hai phần ( $z_{1}, z_{2}$ ), trong đó $z_{1}=\sigma_{1} \sigma_{2} \cdots \sigma_{p(n)}$ đề cập đến vùng bên trái của ô bắt đầu, không bao gồm ô bắt đầu, và $z_{2}=\sigma_{p(n)+1} \sigma_{p(n)+2} \cdots \sigma_{2 p(n)+1}$ đề cập đến phần còn lại của vùng băng thiết yếu. Điều này cung cấp cho chúng ta một cấu hình đã được sửa đổi dưới dạng $\left(q, h, z_{1} z_{2}\right)$. Vì một lý do thực tiễn, ta tiếp tục thay đổi nó thành $\left(z_{2}, z_{1}, h, q\right)$, mà ta gọi là cấu hình lệch. Liên quan đến sự thay đổi này của các cấu hình, ta cũng sửa đổi toán tử tiến hóa thời gian ban đầu, $U_{\delta}$, của $M$ để nó hoạt động trên các cấu hình lệch. Cần lưu ý rằng toán tử mới này không thể được hiện thực hóa bởi mô hình QTM chuẩn. Để phân biệt nó với toán tử đánh giá thời gian ban đầu $U_{\delta}$ của $M$, ta gọi nó là toán tử tiến hóa thời gian lệch và viết nó là $\widehat{U}_{\delta}$.

Bổ đề 4.2 Cho $f$ là bất kỳ hàm lượng tử nào trong $\square_{1}^{\mathrm{QP}}$. Tồn tại một đa thức $p$ và một QTM băng đơn, bảo thủ, thuần túy $M$ tạo ra đầu ra sạch với biên độ $\tilde{\mathbb{C}}$ sao cho, đối với bất kỳ trạng thái lượng tử $|\phi\rangle$ nào trong $\mathcal{H}_{2^{n}}, M$ bắt đầu với một trạng thái lượng tử khác rỗng $|\phi\rangle$ được đưa ra trên băng đầu vào/công việc của nó và, khi nó dừng lại sau $p(n)$ bước, chồng chất $\left|\eta_{M, \phi}\right\rangle$ của các cấu hình lệch cuối cùng của $M$ trên $|\phi\rangle$ có dạng $f(|\phi\rangle) \otimes\left|0, q_{f}\right\rangle$, trong đó $f(|\phi\rangle)$ xuất hiện như nội dung của băng từ ô bắt đầu về bên phải cho đến ký hiệu trống đầu tiên và $q_{f}$ là trạng thái nội duy nhất của $M$.

Chứng minh. Trước tiên, chúng ta chứng minh rằng tất cả các hàm ban đầu trong Sơ đồ I có thể được tính chính xác trong thời gian đa thức trên các QTM băng đơn, có biên độ $\tilde{\mathbb{C}}$, bảo thủ, thuần túy có đầu ra sạch trên bảng chữ cái đầu vào/đầu ra $\{0,1\}$. Để rõ ràng, mục tiêu của chúng ta ở đây là chứng minh rằng, đối với mỗi hàm lượng tử $f$ trong Sơ đồ I, tồn tại một QTM với các tính chất được nêu trong bổ đề sao cho một chồng chất $\left|\eta_{M, \phi}\right\rangle$ của các cấu hình lệch cuối cùng của $M$ là dạng $f(|\phi\rangle) \otimes\left|0, q_{f}\right\rangle$, trong đó $f(|\phi\rangle)$ là nội dung của băng $M$ từ ô bắt đầu về bên phải cho đến ký hiệu trống đầu tiên.

Vì $I$ (đồng nhất) dễ mô phỏng, hãy xem xét $P H A S E_{\theta}$. Trong trường hợp này, ta lấy một QTM áp dụng một chuyển đổi có dạng $\delta\left(q_{0}, \sigma\right)=e^{i \theta \sigma}\left|q_{f}, \sigma, N\right\rangle$ cho bất kỳ bit $\sigma \in\{0,1\}$. Rõ ràng, nếu ta bắt đầu với một cấu hình lệch ban đầu $|\phi\rangle|0\rangle\left|q_{0}\right\rangle$, thì ta dừng lại với $P H A S E_{\theta}(|\phi\rangle) \otimes|0\rangle\left|q_{f}\right\rangle$; nói ngắn gọn, QTM này "tính chính xác" $P H A S E_{\theta}$. Đối với $R O T_{\theta}$, ta sử dụng chuyển đổi được định nghĩa bởi $\delta\left(q_{0}, \sigma\right)=\cos \theta\left|q_{f}, \sigma, N\right\rangle+(-1)^{\sigma} \sin \theta\left|q_{f}, \bar{\sigma}, N\right\rangle$ để tính chính xác $R O T_{\theta}$. Để mô phỏng $N O T$, ta định nghĩa một QTM có chuyển đổi $\delta\left(q_{0}, \sigma\right)=\left|q_{f}, \bar{\sigma}, N\right\rangle$. Trong trường hợp $S W A P$, ta cần chuẩn bị các trạng thái nội $p_{\sigma_{1}}$ và $r_{\sigma_{2}}$ cũng như các chuyển đổi sau: $\delta\left(q_{0}, \sigma_{1}\right)=\left|p_{\sigma_{1}}, \#, R\right\rangle, \delta\left(p_{\sigma_{1}}, \sigma_{2}\right)=\left|r_{\sigma_{2}}, \sigma_{1}, L\right\rangle, \delta\left(r_{\sigma_{2}}, \#\right)=\left|q_{f}, \sigma_{2}, N\right\rangle$ cho các bit $\sigma_{1}, \sigma_{2} \in\{0,1\}$. Đối với $M E A S[a](|\phi\rangle)$, ta bắt đầu bằng cách kiểm tra qubit đầu tiên của $|\phi\rangle$. Nếu nó không phải là $a$, thì ta cho QTM từ chối đầu vào; ngược lại, ta không làm gì cả. Một cách chính thức hơn, ta định nghĩa $\delta\left(q_{0}, \bar{a}\right)=\left|q_{\text {rej }}, \bar{a}, N\right\rangle$ và $\delta\left(q_{0}, a\right)=\left|q_{f}, a, N\right\rangle$. Không khó để thấy rằng $\left|\eta_{M, \phi}\right\rangle$ có dạng $f(|\phi\rangle) \otimes\left|0, q_{0}\right\rangle$.

Tiếp theo, chúng tôi muốn mô phỏng từng quy tắc xây dựng trên một QTM thích hợp. Theo giả thuyết quy nạp, tồn tại ba QTM băng đơn, thời gian đa thức, có biên độ $\tilde{\mathbb{C}}$, bảo thủ, thuần túy $M_{g}, M_{h}$, và $M_{p}$ thỏa mãn bổ đề cho $g, h$, và $p$, tương ứng. Bằng cách cài đặt một đồng hồ nội bộ theo cách nào đó với một đa thức $r$ nhất định, chúng tôi có thể làm cho $M_{g}, M_{h}$, và $M_{p}$ dừng lại chính xác trong $r(n)$ thời gian trên bất kỳ đầu vào nào có độ dài $n \in \mathbb{N}$. Trong phần tiếp theo, hãy để $|\phi\rangle$ biểu thị bất kỳ đầu vào nào trong $\mathcal{H}_{\infty}$.\\[0pt]
[composition] Xem xét trường hợp $f=C o m p o[g, h]$. Chúng tôi tính $f$ như sau. Trước tiên, chúng tôi chạy $M_{h}$ trên $|\phi\rangle$ và nhận được một chồng chất các cấu hình lệch cuối cùng, ví dụ, $h(|\phi\rangle) \otimes\left|0, q_{f}^{\prime}\right\rangle$ cho một trạng thái nội duy nhất $q_{f}^{\prime}$ của $M_{h}$. Vì $M_{h}$ là tĩnh và ở dạng chuẩn, chúng tôi có thể tiếp tục chạy $M_{g}$ trên trạng thái lượng tử kết quả\\
$h(|\phi\rangle)$, coi $q_{f}^{\prime}$ là trạng thái nội ban đầu mới và tạo ra $\left|\eta_{M, \phi}\right\rangle=g(h(|\phi\rangle)) \otimes\left|0, q_{f}\right\rangle$ cho một trạng thái nội duy nhất $q_{f}$ của $M_{g}$.\\[0pt]
[branching] Giả sử rằng $f=\operatorname{Branch}[g, h]$. Chúng tôi kiểm tra qubit đầu tiên của $|\phi\rangle$. Nếu nó là 0 , thì chúng tôi chạy $M_{g}$ trên $\langle 0 \mid \phi\rangle$; ngược lại, chúng tôi chạy $M_{h}$ trên $\langle 1 \mid \phi\rangle$. Điều này tạo ra $|0\rangle \otimes g(\langle 0 \mid \phi\rangle) \otimes\left|0, q_{h, f}\right\rangle+|1\rangle \otimes h(\langle 1 \mid \phi\rangle) \otimes\left|0, q_{g, f}\right\rangle$, trong đó $q_{g, f}$ và $q_{h, f}$ lần lượt là các trạng thái nội duy nhất của $M_{g}$ và $M_{h}$. Lưu ý rằng chúng tôi không cần kiểm tra xem $\ell(|\phi\rangle) \geq 2$ hay không. Một cách chính thức, chúng tôi thêm các chuyển đổi sau vào các chuyển đổi của $M_{g}$ và $M_{h}: \delta\left(q_{0}, 0\right)=\left|q_{0}, \#_{g}, R\right\rangle$ và $\delta\left(q_{0}, 1\right)=\left|q_{0}, \#_{h}, R\right\rangle$, trong đó $\#_{g}$ và $\#_{h}$ là các ký hiệu được chỉ định đánh dấu bên trái của các "ô bắt đầu" mới cho $M_{g}$ và $M_{h}$, tương ứng. Khi kết thúc, đầu băng di chuyển trở lại ô bắt đầu. Sau đó, chúng tôi thay thế $\#_{g}$ và $\#_{h}$ lần lượt bằng 0 và 1 bằng cách đi vào một trạng thái nội mới $q_{f}$; cụ thể, $\delta\left(q_{g, f}, \#_{g}\right)=\left|q_{f}, 0, N\right\rangle$ và $\delta\left(q_{h, f}, \#_{h}\right)=\left|q_{f}, 1, N\right\rangle$. Việc chuyển đổi này là có thể một cách đơn vị.\\[0pt]
[multi-qubit quantum recursion] Cuối cùng, hãy chứng minh một mô phỏng của phép đệ quy lượng tử đa qubit được giới thiệu bởi $f=k Q R e c_{t}\left[g, h, p \mid \mathcal{F}_{k}\right]$ với $\mathcal{F}_{k}=\left\{f_{s}\right\}_{s \in\{0,1\}^{k}}$. Để dễ đọc, hãy xem xét trường hợp đơn giản nhất khi $k=1, f_{0}=f$, và $f_{1}=I$. Các trường hợp khác có thể được xử lý tương tự. Xem xét QTM "nhiều băng" $M_{f}$ hoạt động đại khái như mô tả dưới đây. Hãy để $T$ là một hằng số dương đủ lớn.\\
(1) Trong giai đoạn ban đầu này, bắt đầu với đầu vào $|\phi\rangle$, chúng tôi chuẩn bị một bộ đếm trên một băng công việc mới và một đồng hồ nội bộ trên một băng công việc khác. Chúng tôi sử dụng đồng hồ để điều chỉnh thời điểm kết thúc của tất cả các đường tính toán. Đếm số lượng $\ell(|\phi\rangle)$ các qubit đơn giản bằng cách tăng bộ đếm khi di chuyển đầu băng đầu vào từ ô bắt đầu sang bên phải. Ban đầu, chúng tôi đặt trạng thái lượng tử hiện tại, ví dụ, $|\xi\rangle$ biểu thị trên băng là $|\phi\rangle$ và đặt bộ đếm hiện tại $k$ là $\ell(|\xi\rangle)(=n)$. Đi đến giai đoạn phân chia.\\
(2) Trong giai đoạn phân chia này, chúng tôi thực hiện theo cách quy nạp thủ tục sau bằng cách sử dụng đồng hồ. (\textit{) Giả sử rằng băng đầu vào hiện tại chứa một trạng thái lượng tử $|\xi\rangle$ và bộ đếm có $k=\ell(|\xi\rangle)$. Nếu $k \leq t$, thì chờ cho đến khi đồng hồ đạt $T$ và sau đó đi đến giai đoạn xử lý. Ngược lại, chạy $M_{p}$ trên $|\xi\rangle$ để tạo ra $\left|\psi_{p, \xi}\right\rangle$ và quan sát qubit đầu tiên của $\left|\psi_{p, \xi}\right\rangle$ trong cơ sở tính toán, thu được $\left\langle b \mid \psi_{p, \xi}\right\rangle$ cho mỗi $b \in\{0,1\}$. Nếu $b$ là 1 , thì chạy $M_{h}$ trên $|1\rangle\left\langle 1 \mid \psi_{p, \xi}\right\rangle$, thu được $h\left(|1\rangle\left\langle 1 \mid \psi_{p, \xi}\right\rangle\right)$, được coi là $f(|1\rangle\langle 1 \mid \xi\rangle)$, chờ cho đến khi đồng hồ đạt $T$, và sau đó bắt đầu giai đoạn xử lý. Ngược lại, khi $b$ là 0 , di chuyển bit 0 này sang một băng riêng biệt để ghi nhớ và sau đó cập nhật cả $|\xi\rangle$ và $k$ thành $\left\langle 0 \mid \psi_{p, \xi}\right\rangle$ và $k-1$, tương ứng. Tiếp tục (}).\\
(3) Trong giai đoạn xử lý này, chúng tôi bắt đầu với một trạng thái lượng tử $|\xi\rangle$, được tạo ra trong giai đoạn phân chia. Cho $k=\ell(|\xi\rangle)$. Chúng tôi thực hiện theo cách quy nạp thủ tục sau. (\textbf{) Nếu $k \leq t$, thì chúng tôi chạy $M_{g}$ trên đầu vào $|\xi\rangle$ và tạo ra $g(|\xi\rangle)$, được coi là $f(|\xi\rangle)$. Cập nhật $|\xi\rangle$ thành trạng thái lượng tử kết quả. Ngược lại, chúng tôi di chuyển lại bit 0 cuối cùng đã được lưu trữ từ băng riêng biệt, chạy $M_{h}$ trên $|0\rangle|\xi\rangle$, thu được $h(|0\rangle|\xi\rangle)$, và sau đó cập nhật $|\xi\rangle$ thành trạng thái lượng tử đã thu được và $k$ thành $k+1$ vì $\ell(|0\rangle|\xi\rangle)=k+1$. Nếu tất cả các $b$ đã được tiêu thụ (tương đương, $k=n$ ), thì chờ cho đến khi đồng hồ đạt $2 T$, xuất $|\xi\rangle$, và dừng. Ngược lại, tiếp tục (}).

Thời gian chạy của QTM trên được giới hạn trên bởi một đa thức nhất định theo độ dài $\ell(|\phi\rangle)$ vì mỗi thủ tục (*) và (**) được lặp lại nhiều nhất $\ell(|\phi\rangle)$ lần. Mặc dù $M_{f}$ lưu trữ các bit trên băng riêng biệt, nhưng tất cả các bit đó đều được di chuyển trở lại và sử dụng hết vào cuối quá trình tính toán. Thực tế này cho thấy một chồng chất các cấu hình lệch cuối cùng của $M_{f}$ có dạng $f(|\phi\rangle) \otimes\left|0, q_{f}\right\rangle$ cho một trạng thái nội duy nhất $q_{f}$. Cuối cùng, chúng tôi chuyển đổi QTM đa băng này thành một QTM băng đơn tương đương về mặt tính toán có các tính chất mong muốn được nêu trong bổ đề.

Điều này hoàn thành chứng minh của Bổ đề 4.2.\\
Vì lý do ký hiệu, chúng tôi viết $M[r]$ để chỉ chuỗi đầu ra được viết trong một cấu hình lệch cuối cùng $r$, bao gồm chỉ vùng băng thiết yếu của $M$. Chúng tôi viết $F S C_{M, n}$ để biểu thị tập hợp tất cả các cấu hình lệch cuối cùng có thể có của $M$ được tạo ra cho bất kỳ đầu vào nào có độ dài $n$.

Chìa khóa để chứng minh Định lý 4.1 là bổ đề sau đây, đảm bảo sự tồn tại của một hàm $\widehat{\square_{1}^{\mathrm{QP}}}$ $g$ mà kết quả của nó $\left.\left.g\left(\mid \phi^{p}(|\psi\rangle)\right\rangle\right)\left(=\mid \phi_{g}^{p}(|\psi\rangle)\right\rangle\right)$ "gần như" đặc trưng cho cấu hình lệch cuối cùng đã mã hóa $\sum_{x:|x|=n} \sum_{r \in F S C_{M, n}}\langle x \mid \psi\rangle|\widetilde{M[r]}\rangle\left|\xi_{x, r}\right\rangle$ của $M$, trong đó $\widetilde{M[r]}$ là mã hóa của $M[r]$. Ở đây, mã hóa $\widetilde{M[r]}$ là cần thiết để xây dựng $g$ trong chứng minh của bổ đề. Tuy nhiên, như cũng được chỉ ra trong chứng minh, việc thực hiện một thủ tục giải mã bổ sung cho $\widetilde{M[r]}$ cho phép chúng ta thay thế $\widetilde{M[r]}$ bằng $M[r]$.

Bổ đề 4.3 (Bổ đề chính) Cho $M$ là một QTM băng đơn, thời gian đa thức, có biên độ $\tilde{\mathbb{C}}$, bảo thủ, thuần túy có đầu ra sạch trên bảng chữ cái đầu vào/đầu ra $\Sigma=\{0,1\}$ và bảng chữ cái băng $\Gamma=\{0,1, b\}$. Giả sử rằng, khi $M$ dừng lại trên đầu vào $|\psi\rangle$ có độ dài $n$, một chồng chất các cấu hình lệch cuối cùng đã mã hóa có dạng\\
$\sum_{x:|x|=n} \sum_{r \in F S C_{M, n}}\langle x \mid \psi\rangle|\widetilde{M[r]}\rangle\left|\xi_{x, r}\right\rangle$, trong đó $\left|\xi_{x, r}\right\rangle$ biểu thị một trạng thái lượng tử thích hợp mô tả phần còn lại của các cấu hình lệch cuối cùng đã mã hóa ngoài $\widetilde{M[r]}$. Tồn tại một hàm lượng tử $g$ trong $\widehat{\square_{1}^{\mathrm{QP}}}$ và một đa thức $p$ sao cho, đối với bất kỳ số $n \in \mathbb{N}^{+}$và mọi trạng thái lượng tử $\left.|\psi\rangle \in \mathcal{H}_{2^{n}}, \mid \phi_{g}^{p}(|\psi\rangle)\right\rangle$ có dạng $\sum_{x:|x|=n} \sum_{r \in F S C_{M, n}}\langle x \mid \psi\rangle|\widetilde{M[r]}\rangle\left|\widehat{\xi}_{x, r}\right\rangle$ cho các trạng thái lượng tử nhất định $\left\{\left|\widehat{\xi}_{x, r}\right\rangle\right\}_{x, r}$ thỏa mãn rằng, đối với bất kỳ $x, x^{\prime} \in\{0,1\}^{n}$ và $r, r^{\prime} \in F S C_{M, n}$, (i) $\left\|\left\langle\xi_{x, r} \mid \xi_{x^{\prime}, r^{\prime}}\right\rangle\right\|=\left\|\left\langle\widehat{\xi}_{x, r} \mid \widehat{\xi}_{x^{\prime}, r^{\prime}}\right\rangle\right\|$ và (ii) $\left\langle\widehat{\xi}_{x, r} \mid \widehat{\xi}_{x, r^{\prime}}\right\rangle=0$ nếu $r \neq r^{\prime}$. Hơn nữa, có thể sửa đổi $g$ thành $g^{\prime}$, thỏa mãn $\left.\mid \phi_{g^{\prime}}^{p}(|\phi\rangle)\right\rangle=\sum_{x:|x|=n} \sum_{r \in F S C_{M, n}}\langle x \mid \psi\rangle|M[r]\rangle\left|\widehat{\xi}_{x, r}^{\prime}\right\rangle$ cho các trạng thái lượng tử thích hợp $\left\{\left|\widehat{\xi}_{x^{\prime}, r^{\prime}}^{\prime}\right\rangle\right\}_{x^{\prime}, r^{\prime}}$ thỏa mãn các điều kiện (i)-(ii).

Chứng minh của Bổ đề 4.3 khá dài và được hoãn lại cho đến Mục 5 . Trong lúc đó, chúng tôi quay lại Định lý 4.1 và trình bày chứng minh của nó bằng cách sử dụng các Bổ đề 4.2 và 4.3.

Chứng minh Định lý 4.1. Cho $\varepsilon \in[0,1 / 2)$ là bất kỳ hằng số nào và cho $f$ là bất kỳ hàm có giới hạn đa thức nào ánh xạ $\{0,1\}^{*}$ tới $\{0,1\}^{*}$.\\
$(1 \Rightarrow 2)$ Giả sử rằng $f$ thuộc FBQP. Lấy một QTM đa băng, thời gian đa thức, có biên độ $\tilde{\mathbb{C}}$, được hình thành tốt $N$ tính toán $f$ với xác suất lỗi bị chặn. Hãy chọn một đa thức $p$ giới hạn thời gian chạy của $N$ trên mọi đầu vào. Có thể "mô phỏng" $N$ với xác suất thành công cao bằng một QTM băng đơn nhất định, có biên độ $\tilde{\mathbb{C}}$, bảo thủ, thuần túy, có đầu ra sạch sao cho máy lấy đầu vào $x$ và kết thúc bằng việc tạo ra $f(x) b w$ trong vùng bên phải của ô bắt đầu của băng đầu vào/công việc với xác suất lỗi bị chặn, trong đó $b$ là một ký hiệu băng trống duy nhất. Vì bất kỳ QTM lỗi bị chặn nào cũng có thể khuếch đại tự do xác suất thành công của nó, chúng ta có thể giả sử mà không mất tính tổng quát rằng xác suất lỗi của $M$ là nhiều nhất $\varepsilon$. Lưu ý rằng cấu hình lệch cuối cùng đã mã hóa bắt đầu bằng một chuỗi đầu ra. Do đó, chồng chất các cấu hình lệch cuối cùng đã mã hóa của $M$ trên đầu vào $x$ có độ dài $n$ phải có dạng $\sum_{r \in F S C_{M, n}}|\widetilde{M[r]}\rangle\left|\xi_{x, r}\right\rangle$ cho các trạng thái lượng tử được chọn phù hợp $\left\{\left|\xi_{x, r}\right\rangle\right\}_{r \in F S C_{M, n}}$. Vì $M$ tính toán $f$ với xác suất lỗi nhiều nhất $\varepsilon$, chúng ta kết luận rằng $\sum_{r \in F S C_{M, n}} \|\langle\widetilde{f(x)} \mid \widetilde{M[r]}\rangle\left|\xi_{x, r}\right\rangle \|^{2} \geq 1-\varepsilon$ cho tất cả $x$. Bổ đề 4.3 tiếp tục cung cấp cho chúng ta một hàm lượng tử đặc biệt $g \in \widehat{\square_{1}^{\mathrm{QP}}}$ sao cho, đối với bất kỳ số $n \in \mathbb{N}$ và bất kỳ chuỗi $x \in\{0,1\}^{n},\left|\phi_{g}^{p}(x)\right\rangle$ có dạng $\sum_{r \in F S C_{M, n}}|M[r]\rangle\left|\widehat{\xi}_{x, r}\right\rangle$ cho các trạng thái lượng tử nhất định $\left\{\left|\widehat{\xi}_{x, r}\right\rangle\right\}_{r \in F S C_{M, n}}$ thỏa mãn rằng, đối với tất cả $r, r^{\prime} \in F S C_{M, n},\left\|\left\langle\xi_{x, r} \mid \xi_{x, r^{\prime}}\right\rangle\right\|=\left\|\left\langle\widehat{\xi}_{x, r} \mid \widehat{\xi}_{x, r^{\prime}}\right\rangle\right\|$ và $\left\langle\widehat{\xi}_{x, r} \mid \widehat{\xi}_{x, r^{\prime}}\right\rangle=0$ nếu $r \neq r^{\prime}$. Do đó, chúng ta kết luận rằng $\left\|\left\langle f(x) \mid \phi_{g}^{p}(x)\right\rangle\right\|^{2}=\sum_{r \in F S C_{M, n}} \|\langle f(x) \mid M[r]\rangle\left|\widehat{\xi}_{x, r}\right\rangle\left\|^{2}=\sum_{r \in F S C_{M, n}}|\langle\widetilde{f(x)} \mid \widetilde{M[r]}\rangle|^{2}\right\|\left|\widehat{\xi}_{x, r}\right\rangle \|^{2}$ vì và $\langle f(x) \mid M[r]\rangle=\langle\widetilde{f(x)} \mid \widetilde{M[r]}\rangle$ và $\left\langle\widehat{\xi}_{x, r} \mid \widehat{\xi}_{x, r^{\prime}}\right\rangle=0$ nếu $r \neq r^{\prime}$. Hơn nữa, từ $\|\left|\xi_{x, r}\right\rangle\|=\|\left|\widehat{\xi}_{x, r}\right\rangle \|$, ta có được số hạng $|\langle\widetilde{f(x)} \mid \widetilde{M[r]}\rangle|^{2} \|\left|\widehat{\xi}_{x, r}\right\rangle \|^{2}$ bằng $\|\langle\widetilde{f(x)} \mid \widetilde{M[r]}\rangle\left|\xi_{x, r}\right\rangle \|^{2}$. Do đó, $\left\|\left\langle f(x) \mid \phi_{g}^{p}(x)\right\rangle\right\|^{2}$ ít nhất là $1-\varepsilon$.\\
( $2 \Rightarrow 3$ ) Cho $\varepsilon$ là bất kỳ hằng số nào trong $[0,1 / 2)$ và đặt $\varepsilon^{\prime}=1-\sqrt{1-\varepsilon}$. Hãy chọn một đa thức $p$ và một hàm $g \in \widehat{\square_{1}^{\mathrm{QP}}}$ sao cho $|f(x)| \leq p(|x|)$ và $\left\|\left\langle f(x) \mid \phi_{g}^{p}(x)\right\rangle\right\|^{2} \geq 1-\varepsilon^{\prime}$ cho tất cả các chuỗi $x \in\{0,1\}^{*}$. Bắt đầu với đầu vào $\left.\left|\phi^{p, f}(x)\right\rangle\left(=\left|\widetilde{\left.0^{|f(x)|}\right\rangle}\right| 0^{|f(x)|+1} 1\right\rangle\left|\phi^{p}(x)\right\rangle\right)$, trước tiên chúng ta áp dụng $g$ vào phần cuối $\left|\phi^{p}(x)\right\rangle$ của $\left|\phi^{p, f}(x)\right\rangle$ và thu được $\left|\widetilde{0^{|f(x)|}}\right\rangle\left|0^{|f(x)|+1} 1\right\rangle\left|\phi_{g}^{p}(x)\right\rangle$. Tiếp theo, chúng ta áp dụng Bổ đề 5.1 để mã hóa $|f(x)|$ qubit đầu tiên của $\left|\phi_{g}^{p}(x)\right\rangle$ với sự trợ giúp của $\left|0^{|f(x)|+1} 1\right\rangle$ và sau đó thu được $\left|\eta_{f, x}\right\rangle=\sum_{s:|s|=|f(x)|}\left|\widetilde{0^{|f(x)|}}\right\rangle|\tilde{s}\rangle \otimes\left\langle s \mid \phi_{g}^{p}(x)\right\rangle$. Sử dụng hàm lượng tử $C O P Y_{2}$ được đưa ra trong Bổ đề 5.2, chúng ta sau đó sao chép mỗi qustring $|\tilde{s}\rangle$ của $\left|\eta_{f, x}\right\rangle$ lên $\left|\widetilde{0^{|f(x)|}}\right\rangle$ và tạo ra một trạng thái lượng tử có dạng $\sum_{s:|s|=|f(x)|}\left(|\tilde{s}\rangle|\tilde{s}\rangle \otimes\left\langle s \mid \phi_{g}^{p}(x)\right\rangle\right)$. Sau đó, chúng ta giải mã lần xuất hiện đầu tiên và lần xuất hiện thứ hai của $|\tilde{s}\rangle$ thành $\left|0^{|f(x)|+1} 1\right\rangle|s\rangle$. Vì tiện lợi, chúng ta biến đổi lần xuất hiện đầu tiên của $\left|0^{|f(x)|+1} 1\right\rangle|s\rangle$ thành $|s\rangle\left|0^{|f(x)|+1} 1\right\rangle$. Trong phần tiếp theo, chúng ta viết tắt $|s\rangle\left\langle s \mid \phi_{g}^{p}(x)\right\rangle$ là $\left|\zeta_{g}^{p}[s]\right\rangle$. Cuối cùng, chúng ta áp dụng cục bộ $g^{-1}$ vào phần cuối $\left|\zeta_{g}^{p}[s]\right\rangle$, tạo ra $\left|\xi_{x}\right\rangle=\sum_{s:|s|=|f(x)|}\left(|s\rangle\left|\left(0^{|f(x)|+1} 1\right)^{2}\right\rangle \otimes g^{-1}\left(\left|\zeta_{g}^{p}[s]\right\rangle\right)\right)$.

Cho $\left|\Psi_{f(x)}\right\rangle=|f(x)\rangle\left|\left(0^{|f(x)|+1} 1\right)^{2}\right\rangle\left|\phi^{p}(x)\right\rangle$. Vì $\left|\phi_{g}^{p}(x)\right\rangle=\sum_{s:|s|=|f(x)|}\left|\zeta_{g}^{p}[s]\right\rangle$ và $g^{-1}\left(\left|\phi_{g}^{p}(x)\right\rangle\right)= \left|\phi^{p}(x)\right\rangle$, ta suy ra $\left|\phi^{p}(x)\right\rangle=\sum_{s:|s|=|f(x)|} g^{-1}\left(\left|\zeta_{g}^{p}[s]\right\rangle\right)$. Do đó, $\left|\Psi_{f(x)}\right\rangle$ trùng khớp với $\sum_{s:|s|=|f(x)|}\left(|f(x)\rangle\left|\left(0^{|f(x)|+1} 1\right)^{2}\right\rangle \otimes g^{-1}\left(\left|\zeta_{g}^{p}[s]\right\rangle\right)\right)$. Hãy xem xét tích vô hướng $\left\langle\Psi_{f(x)} \mid \xi_{x}\right\rangle$. Bằng một phép tính đơn giản, $\left\langle\Psi_{f(x)} \mid \xi_{x}\right\rangle$ bằng $\sum_{s:|s|=|f(x)|}\langle f(x) \mid s\rangle \otimes \tau_{x}(s)$, trong đó $\tau_{x}(s)$ là tích vô hướng giữa $\left|\phi^{p}(x)\right\rangle$ và $g^{-1}\left(\left|\zeta_{g}^{p}[s]\right\rangle\right)$. Vì $g^{-1}$ thuộc về $\widehat{\square_{1}^{\mathrm{QP}}}$ và bảo toàn chiều cũng như bảo toàn chuẩn theo Mệnh đề 3.5 và Bổ đề 3.4, $\tau_{x}(s)$ bằng $\left\langle\zeta_{g}^{p}[s] \mid \zeta_{g}^{p}[s]\right\rangle$. Vì $s$ bị ép phải lấy giá trị $f(x)$ trong $\left\langle\Psi_{f(x)} \mid \xi_{x}\right\rangle$, ta có được $\left\langle\Psi_{f(x)} \mid \xi_{x}\right\rangle=\left\langle\zeta_{g}^{p}[f(x)] \mid \zeta_{g}^{p}[f(x)]\right\rangle$, bằng $\left\|\left\langle f(x) \mid \phi_{g}^{p}(x)\right\rangle\right\|^{2}$. Do đó, ta kết luận rằng $\left|\left\langle\Psi_{f(x)} \mid \xi_{x}\right\rangle\right|^{2}=\left\|\left\langle f(x) \mid \phi_{g}^{p}(x)\right\rangle\right\|^{4} \geq\left(1-\varepsilon^{\prime}\right)^{2}=1-\varepsilon$, như yêu cầu.\\
( $3 \Rightarrow 1$ ) Vì $f$ có giới hạn đa thức, hãy lấy một đa thức $p$ sao cho $|f(x)| \leq p(|x|)$ đúng với tất cả các chuỗi\\
$x$. Vì chúng ta không biết độ dài của $f(x)$, chúng ta muốn mở rộng $f$ bằng cách đặt $f_{1}(x)=f(x) 01^{p(|x|)+2-|f(x)|}$ để $\left|f_{1}(x)\right|=p(|x|)+2$ cho tất cả các chuỗi $x$. Cố định $\varepsilon \in[0,1 / 2)$. Hãy giả sử rằng tồn tại một hàm $g_{1} \in \widehat{\square_{1}^{\mathrm{QP}}}$ và một đa thức $p_{1}$ thỏa mãn $\left|f_{1}(x)\right| \leq p_{1}(|x|)$ và $\left|\left\langle\Psi_{f_{1}(x)} \mid \phi_{g}^{p_{1}, f_{1}}(x)\right\rangle\right|^{2} \geq 1-\varepsilon$ cho tất cả các trường hợp $x \in\{0,1\}^{*}$.

Sử dụng Bổ đề 4.2 cho hàm lượng tử $g_{1}$, chúng ta có thể lấy một QTM băng đơn, thời gian đa thức, bảo thủ, thuần túy có biên độ $\tilde{\mathbb{C}}$ cho $M$ mà $M$ trên đầu vào $|\phi\rangle$ tạo ra đầu ra sạch của $g_{1}(|\phi\rangle)$ trên băng của nó. Chúng ta xem xét máy sau. Trên đầu vào $x \in \Sigma^{*}$, trước tiên chúng ta tính giá trị $p(|x|)$ một cách xác định và tạo ra $\left|\phi^{p, f_{1}}\right\rangle=\left|0^{p(|x|)+2} 1\right\rangle\left|0^{p(|x|)} 10^{9 p(|x|)} 1\right\rangle \otimes|x\rangle$. Sau đó, chúng ta chạy $M$ trên $\left|\phi^{p, f_{1}}(x)\right\rangle$ để tạo ra $\left|\phi_{g_{1}}^{p, f_{1}}(x)\right\rangle$. Vì $\left|\left\langle\Psi_{f_{1}(x)} \mid \phi_{g_{1}}^{p, f_{1}}(x)\right\rangle\right|^{2} \geq 1-\varepsilon,\left|\phi_{g_{1}}^{p, f_{1}}(x)\right\rangle$ chứa $f_{1}(x)$ với xác suất ít nhất $1-\varepsilon$. Từ $f_{1}(x)$, chúng ta trích xuất $f(x)$ và xuất nó. Điều này kết luận rằng $f$ thuộc $F B Q P$.

\subsection*{4.2 Định lý dạng chuẩn lượng tử}
Các bổ đề chính của chúng tôi, Bổ đề 4.2 và 4.3, tiếp tục dẫn đến một phiên bản lượng tử của định lý dạng chuẩn của Kleene 19, 20, khẳng định sự tồn tại của một vị từ đệ quy nguyên thủy $T(e, x, y)$ và một hàm đệ quy nguyên thủy $U(y)$ sao cho, đối với bất kỳ hàm đệ quy $f(x)$ nào, một chỉ số thích hợp (gọi là số Gödel) $e \in \mathbb{N}$ thỏa mãn $f(x)=U(\mu y . T(e, x, y))$ cho tất cả các đầu vào $x \in \mathbb{N}$, trong đó $\mu$ là toán tử tối thiểu hóa. Phát biểu này về cơ bản tương đương với sự tồn tại của máy Turing phổ quát 32. Ở đây, chúng tôi muốn chứng minh một dạng yếu hơn của định lý dạng chuẩn lượng tử bằng cách sử dụng các Bổ đề 4.2, 4.3.\\
Định lý 4.4 (Định lý dạng chuẩn lượng tử) Tồn tại một hàm lượng tử $f$ trong $\square_{1}^{\mathrm{QP}}$ sao cho, đối với bất kỳ hàm lượng tử $g$ nào trong $\square_{1}^{\mathrm{QP}}$ và bất kỳ hằng số $\varepsilon \in(0,1 / 2)$ nào, tồn tại một chuỗi nhị phân $e$ và một đa thức $p$ thỏa mãn $\|\left|\psi_{g, x}\right\rangle\left\langle\psi_{g, x}\right|-\operatorname{tr}_{n}\left(\left|\eta_{f, x}\right\rangle\left\langle\eta_{f, x}\right|\right) \|_{\mathrm{tr}} \leq \varepsilon$ cho bất kỳ đầu vào nào $|x\rangle$ với $x \in\{0,1\}^{n}$, trong đó $\left|\psi_{g, x}\right\rangle=g(|x\rangle)$ và $\left|\eta_{f, x}\right\rangle=f\left(|\tilde{e}\rangle\left|0^{p(|x|)} 1\right\rangle|x\rangle\right)$. Hàm như vậy được gọi là phổ quát.

Thuật ngữ bổ sung $\left|0^{p(|x|)} 1\right\rangle$ trong $\left|\eta_{f, x}\right\rangle$ là cần thiết để cung cấp cho $g$ đủ không gian làm việc như trong trường hợp của Định lý 4.1. Để chứng minh định lý, chúng tôi sử dụng thực tế là có một QTM phổ quát, có thể mô phỏng tất cả các QTM băng đơn được hình thành tốt với độ trễ đa thức với bất kỳ độ chính xác mong muốn nào. Một máy phổ quát như vậy đã được xây dựng bởi Bernstein và Vazirani [5, Định lý 7.1] và bởi Nishimura và Ozawa [26, Định lý 4.1]. Chúng tôi nói rằng $M_{1}$ trên đầu vào $x_{1}$ mô phỏng $M_{2}$ trên đầu vào $x_{2}$ với độ chính xác nhiều nhất $\varepsilon$ nếu khoảng cách biến thiên toàn phần giữa hai phân bố xác suất $\left\{\left\|\langle y| U_{M_{1}}^{p_{1}\left(\left|x_{1}\right|\right)}\left|c_{0,1}^{\left(x_{1}\right)}\right\rangle\right\|^{2}\right\}_{y \in\{0,1\} \leq n}$ và $\left\{\left\|\langle y| U_{M_{2}}^{p_{2}\left(\left|x_{2}\right|\right)}\left|c_{0,2}^{\left(x_{2}\right)}\right\rangle\right\|^{2}\right\}_{y \in\{0,1\} \leq n}$ là nhiều nhất $\varepsilon$, trong đó $n=\max \left\{p_{1}\left(\left|x_{1}\right|\right), p_{2}\left(\left|x_{2}\right|\right)\right\}$ và, cho mỗi chỉ số $i \in\{1,2\}$, $U_{M_{i}}$ là toán tử tiến hóa thời gian của $M_{i}, p_{i}(\cdot)$ biểu thị thời gian chạy của $M_{i}, c_{0, i}^{\left(x_{i}\right)}$ là cấu hình lệch ban đầu của $M_{i}$ trên đầu vào $x_{i}$, và $y$ trải rộng trên tất cả các chuỗi đầu ra có thể có của $M_{i}$.

Mệnh đề 4.5 [5, 17, 25, 26] Tồn tại một QTM băng đơn, được hình thành tốt, tĩnh sao cho, đối với mỗi hằng số $\varepsilon \in(0,1)$, một số $t \in \mathbb{N}$, một QTM băng đơn được hình thành tốt $M$ có biên độ $\tilde{\mathbb{C}}$, $U$ trên đầu vào $\langle M, x, t, \varepsilon\rangle$ mô phỏng $M$ trên đầu vào $x$ trong $t$ bước với độ chính xác nhiều nhất $\varepsilon$ với độ trễ của một đa thức trong $t$ và $\log (1 / \varepsilon)$, trong đó $\langle M, x, t, \varepsilon\rangle$ đề cập đến một mã hóa cố định, hiệu quả của bộ tứ ( $M, x, t, \varepsilon$ ). QTM như vậy được gọi là phổ quát.

Hệ số cải tiến $\log (1 / \varepsilon)$ trong Mệnh đề 4.5 được gán cho Kitaev 17 và Solovay (trích dẫn trong 25 , Phụ lục 3]).

Đối với chứng minh Định lý 4.4, chúng ta cần mô phỏng một QTM phổ quát $U$ được cung cấp bởi Mệnh đề 4.5 trên một QTM bảo thủ nhất định ngay cả với độ chính xác thấp hơn. Về hình thức đầu vào được đưa cho $f$, chúng ta cần chia bộ tứ ( $M, x, t, \varepsilon$ ) trong mệnh đề thành ba phần ( $M, \varepsilon$ ), $t$, và $x$ và sau đó sửa đổi $U$ để $U$ có thể nhận bất kỳ đầu vào nào có dạng $|\tilde{e}\rangle\left|0^{t} 1\right\rangle|x\rangle$, trong đó $e=\langle M, \varepsilon\rangle$, và bắt chước $M$ trên $x$ trong thời gian $t$ với độ chính xác nhiều nhất $\varepsilon$. Hơn nữa, chúng ta cần buộc $U$ tạo ra đầu ra sạch bằng cách di chuyển tất cả các ký hiệu không trống xuất hiện trong vùng bên trái của bất kỳ chuỗi đầu ra nào đến nơi khác trong thời gian đa thức trong $t$.

Định lý 4.4 hiện nay được suy ra trực tiếp bằng cách kết hợp các Bổ đề 4.24 .3 và Mệnh đề 4.5.\\
Chứng minh Định lý 4.4. Như đã giải thích ở trên, cho $U$ biểu thị một QTM phổ quát đã được sửa đổi có thể nhận các đầu vào có dạng $|\tilde{e}\rangle\left|0^{t} 1\right\rangle|x\rangle$ cho bất kỳ số $e, t \in \mathbb{N}^{+}$và bất kỳ chuỗi $x \in\{0,1\}^{*}$ nào và tạo ra đầu ra sạch. Cho bất kỳ hàm lượng tử $g \in \square_{1}^{\mathrm{QP}}$ nào, Bổ đề 4.2 đảm bảo sự tồn tại của một đa thức $r$ và một QTM băng đơn, có biên độ $\tilde{\mathbb{C}}$, bảo thủ, thuần túy có đầu ra sạch cho $M$ mà, trên bất kỳ đầu vào nào $x \in\{0,1\}^{n}$, $M$ tạo ra trong $r(n)$ bước một chồng chất $g\left(\left|0^{p(n)} 1\right\rangle|x\rangle\right) \otimes\left|0, q_{f}\right\rangle$ của các cấu hình lệch cuối cùng bao gồm\\
của vùng băng thiết yếu của $M$ và trạng thái nội của $M$. Để chính xác hơn, chúng tôi ký hiệu $\widehat{U}_{M}$ là toán tử tiến hóa thời gian lệch của $M$ và $c_{0, M}^{(x)}$ là cấu hình lệch ban đầu của $M$ trên bất kỳ đầu vào nào $|x\rangle$. Chúng tôi tiếp tục đặt $\left|\zeta_{M, x}\right\rangle$ là $\widehat{U}_{M}^{r(n)}\left|c_{0, M}^{(x)}\right\rangle$ và $g\left(\left|0^{p(n)} 1\right\rangle|x\rangle\right)$ là $\left|\psi_{g, x}\right\rangle$. Lưu ý rằng $\left|\zeta_{M, x}\right\rangle$ bằng $\left|\psi_{g, x}\right\rangle \otimes\left|0, q_{f}\right\rangle$. Vì $\left|\psi_{g, x}\right\rangle\left\langle\psi_{g, x}\right|=\operatorname{tr}_{n}\left(\left|\psi_{g, x}\right\rangle\left\langle\psi_{g, x}\right| \otimes\left|0, q_{f}\right\rangle\left\langle 0, q_{f}\right|\right),\left|\psi_{g, x}\right\rangle\left\langle\psi_{g, x}\right|$ có thể được biểu thị là $\operatorname{tr}_{n}\left(\left|\zeta_{M, x}\right\rangle\left\langle\zeta_{M, x}\right|\right)$.

Ngược lại, chúng tôi ký hiệu $\widehat{U}_{\delta}$ là toán tử tiến hóa thời gian lệch của $U$ và $c_{0, U}^{\left(x_{e}\right)}$ là cấu hình lệch ban đầu của $U$ trên đầu vào $\left|x_{e}\right\rangle$ cho $e=\langle M, \varepsilon\rangle$ và $x_{e}=\tilde{e} 0^{r(|x|)} 1 x$. Cho $m=\left|x_{e}\right|$ và đặt $\left|\zeta_{U, x_{e}}\right\rangle$ là $\widehat{U}_{\delta}^{p(m)}\left|c_{0, U}^{\left(x_{e}\right)}\right\rangle$, được viết là $\sum_{r \in F S C_{U, m}}|U[r]\rangle\left|\xi_{x_{e}, r}\right\rangle$ cho một tập hợp thích hợp $\left\{\left|\xi_{x_{e}, r}\right\rangle\right\}_{r}$ của các trạng thái lượng tử trực giao. Mệnh đề 4.5 sau đó đảm bảo rằng, bằng cách lựa chọn một đa thức $s$ thích hợp, khoảng cách biến thiên toàn phần giữa $\left\{\left\|\langle y| \widehat{U}_{M}^{t}\left|c_{0, M}^{(x)}\right\rangle\right\|^{2}\right\}_{y \in\{0,1\}^{n}}$ và $\left\{\left\|\langle y| \widehat{U}_{\delta}^{s(t)}\left|c_{0, U}^{\left(x_{e}\right)}\right\rangle\right\|^{2}\right\}_{y \in\{0,1\}^{n}}$ là nhiều nhất $\varepsilon$ cho bất kỳ số $t \geq 0$ nào.

Chúng tôi áp dụng Bổ đề 4.3 và sau đó nhận được một hàm lượng tử $\square_{1}^{\mathrm{QP}}$ $f$ sao cho, đối với bất kỳ trạng thái lượng tử $|\phi\rangle$ nào, $f(|\phi\rangle)$ biểu thị kết quả của $\widehat{U}_{\delta}^{s(r(n))}$ được áp dụng cho $|\phi\rangle$. Vì tiện lợi, chúng tôi biểu thị $f\left(\left|x_{e}\right\rangle\right)$ là $\left|\eta_{f, x_{e}}\right\rangle$. Bổ đề 4.3 một lần nữa ngụ ý rằng $\left|\eta_{f, x_{e}}\right\rangle=\sum_{r \in F S C_{U, m}}|U[r]\rangle\left|\widehat{\xi}_{x_{e}, r}\right\rangle$ với $\left\langle\xi_{x_{e}, r} \mid \xi_{x_{e}, r^{\prime}}\right\rangle=\left\langle\widehat{\xi}_{x_{e}, r} \mid \widehat{\xi}_{x_{e}, r^{\prime}}\right\rangle$ và $\left\langle\xi_{x, r} \mid \xi_{x, r^{\prime}}\right\rangle=0$ nếu $r \neq r^{\prime}$ cho tất cả $r, r^{\prime} \in F S C_{U, m}$. Do đó, chúng tôi nhận được $\operatorname{tr}_{n}\left(\left|\eta_{f, x_{e}}\right\rangle\left\langle\eta_{f, x_{e}}\right|\right)=\sum_{r, r^{\prime} \in F S C_{U, m}}|U[r]\rangle\left\langle U\left[r^{\prime}\right]\right|$. $\operatorname{tr}\left(\left|\widehat{\xi}_{x_{e}, r}\right\rangle\left\langle\widehat{\xi}_{x_{e}, r^{\prime}}\right|\right)=\sum_{r, r^{\prime} \in F S C_{U, m}}\left\langle\widehat{\xi}_{x_{e}, r^{\prime}} \mid \widehat{\xi}_{x_{e}, r}\right\rangle|U[r]\rangle\left\langle U\left[r^{\prime}\right]\right|$, bằng $\sum_{r \in F S C_{U, m}} \|\left|\widehat{\xi}_{x_{e}, r}\right\rangle \|^{2}|U[r]\rangle\langle U[r]|$. Tương tự, chúng tôi nhận được $\operatorname{tr}_{n}\left(\left|\zeta_{U, x_{e}}\right\rangle\left\langle\zeta_{U, x_{e}}\right|\right)=\sum_{r \in F S C_{U, m}} \|\left|\xi_{x_{e}, r}\right\rangle \|^{2}|U[r]\rangle\langle U[r]|$. Từ những tính toán đó cùng với $\|\left|\xi_{x_{e}, r}\right\rangle\|=\|\left|\widehat{\xi}_{x_{e}, r}\right\rangle \|$, đẳng thức $\operatorname{tr}_{n}\left(\left|\eta_{f, x_{e}}\right\rangle\left\langle\eta_{f, x_{e}}\right|\right)=\operatorname{tr}_{n}\left(\left|\zeta_{U, x_{e}}\right\rangle\left\langle\zeta_{U, x_{e}}\right|\right)$ được suy ra ngay lập tức.

Chúng tôi do đó kết luận rằng $\|\left|\psi_{g, x}\right\rangle\left\langle\psi_{g, x}\right|-\operatorname{tr}_{n}\left(\left|\eta_{f, x}\right\rangle\left\langle\eta_{f, x}\right|\right)\left\|_{\text {tr }}=\right\| \operatorname{tr}_{n}\left(\left|\zeta_{M, x}\right\rangle\left\langle\zeta_{M, x}\right|\right)-\operatorname{tr}_{n}\left(\left|\zeta_{U, x_{e}}\right\rangle\left\langle\zeta_{U, x_{e}}\right|\right) \|_{\text {tr }}$. Số hạng cuối cùng được giới hạn trên bởi $\left.\sum_{y:|y|=n}| |\left|\langle y| \widehat{U}_{\delta}^{s(r(n))}\right| c_{0, U}^{(x)}\right\rangle\left\|^{2}-\right\|\langle y| \widehat{U}_{M}^{r(n)}\left|c_{0, M}^{(x)}\right\rangle \|^{2} \mid$, rõ ràng là nhiều nhất $\varepsilon$. Do đó, $f$ là phổ quát. \(\square\)

\section*{5 Chứng minh Bổ đề chính}
Để hoàn thành chứng minh Định lý 4.1, chúng ta cần chứng minh bổ đề chính, Bổ đề 4.3. Mục này nhằm cung cấp chứng minh mong muốn của bổ đề. Chứng minh của chúng tôi được lấy cảm hứng từ một kết quả của Yao 40, người đã chứng minh một mô phỏng mạch lượng tử của một QTM.

\subsection*{5.1 Mô phỏng chức năng của QTMs}
Cốt lõi của chứng minh Bổ đề 4.3 là một mô phỏng từng bước trực tiếp về hành vi của một QTM băng đơn, có biên độ $\tilde{\mathbb{C}}$, bảo thủ, thuần túy $M=\left(Q, \Sigma, \Gamma, \delta, q_{0}, Q_{f}\right)$, có đầu ra sạch. Để đơn giản, chúng tôi giả sử rằng $\Sigma=\{0,1\}$ và $\Gamma=\{0,1, b\}$, trong đó $b$ ở đây đại diện cho một ký hiệu băng trống duy nhất, thay vì $\#$ được sử dụng trong các phần trước. Chúng tôi cũng giả sử rằng $Q_{f}=\left\{q_{f}\right\}$ và $Q=\{0,1\}^{\ell}$ cho một số chẵn $\ell>0$ nhất định với $q_{0}=0^{\ell}$ và $q_{f}=1^{\ell}$. Hãy giả sử rằng, bắt đầu với chuỗi đầu vào nhị phân $x$ được viết trên băng đầu vào/công việc/đầu ra đơn, $M$ dừng lại trong nhiều nhất $p(|x|)$ bước, trong đó $p$ là một đa thức thích hợp liên quan chỉ với $M$. Chúng tôi cũng giả sử rằng tất cả các đường tính toán của $M$ trên mỗi đầu vào dừng lại đồng thời. Vì tiện lợi, chúng tôi cũng yêu cầu rằng $p(n)>\ell$ cho bất kỳ $n \in \mathbb{N}$ nào. Điều quan trọng là phải nhớ rằng $M$ cuối cùng dừng lại bằng cách đi vào một trạng thái nội duy nhất $q_{f}$ và làm cho đầu băng ở yên tại ô bắt đầu và không có ký hiệu nào không trống xuất hiện trong vùng bên trái của bất kỳ chuỗi đầu ra nào.

Cho $x=x_{1} x_{2} \cdots x_{n}$ là bất kỳ đầu vào nào được đưa cho $M$, trong đó $x_{i}$ là một bit trong $\{0,1\}$ cho mỗi chỉ số $i \in[n]$. Liên quan đến $p$, một vùng băng thiết yếu của $M$ trên $x$ bao gồm tất cả các ô băng được đánh chỉ số giữa $-p(n)$ và $+p(n)$. Chúng tôi biểu thị nội dung băng của vùng băng thiết yếu như vậy là một chuỗi có dạng $\sigma_{1} \sigma_{2} \cdots \sigma_{2 p(n)+1}$ có độ dài chính xác $2 p(|x|)+1$ trên bảng chữ cái băng $\Gamma=\{0,1, b\}$. Chúng tôi theo dõi sự thay đổi của các ký hiệu này khi $M$ thực hiện các bước di chuyển của nó, trong đó $\sigma_{i}$ là một ký hiệu băng được viết tại ô được đánh chỉ số $i-p(n)-1$ cho mỗi chỉ số $i \in[2 p(n)+1]$.

Vì tất cả các hàm $\square_{1}^{\text {QP }}$ được định nghĩa để xử lý các trạng thái lượng tử trong $\mathcal{H}_{\infty}$, chúng tôi cần mã hóa từng ký hiệu băng và do đó là nội dung băng. Chúng tôi do đó định nghĩa một qustring mới mã hóa đúng cấu hình $\gamma=\left(q, h, \sigma_{1} \sigma_{2} \cdots \sigma_{2 p(n)+1}\right)$ của $M$ là

$$
|q\rangle \otimes\left|s_{1}, \hat{\sigma}_{1}\right\rangle \otimes\left|s_{2}, \hat{\sigma}_{2}\right\rangle \otimes\left|s_{3}, \hat{\sigma}_{3}\right\rangle \otimes \cdots \otimes\left|s_{2 p(n)+1}, \hat{\sigma}_{2 p(n)+1}\right\rangle,
$$

trong đó mỗi $\hat{\sigma}_{i}$ nằm trong $\{\hat{0}, \hat{1}, \hat{b}\}$, mỗi $s_{i} \in\{\hat{2}, \hat{3}\}$ biểu thị sự hiện diện của đầu băng (trong đó $\hat{2}$ có nghĩa là "đầu ở đây" và $\hat{3}$ có nghĩa là "không có đầu ở đây") tại ô $i-p(n)-1$, và $q$ là một trạng thái nội trong $Q$. Trong các phần tiếp theo, chúng tôi gọi một qustring như vậy là mã của cấu hình $\gamma$ của $M$ và biểu thị nó là $|\hat{\gamma}\rangle$. Mã $|\hat{\gamma}\rangle$ này có độ dài $\ell(|\hat{\gamma}\rangle)=8 p(n)+\ell+4$, là số chẵn và lớn hơn $n$.

Cho bất kỳ đầu vào nhị phân $x=x_{1} x_{2} \cdots x_{n}$ có độ dài $n$, hãy nhớ lại từ Mục 4.1 rằng $\left|\phi^{p}(x)\right\rangle= \left|0^{|x|} 1\right\rangle\left|0^{p(|x|)} 1\right\rangle\left|0^{11 p(|x|)+6} 1\right\rangle|x\rangle$ và $\left|\phi^{p, f}(x)\right\rangle=\left|\widetilde{0^{|f(x)|}}\right\rangle\left|0^{|f(x)|+1} 1\right\rangle\left|\phi^{p}(x)\right\rangle$. Ngoại trừ Bước 1) trong Mục 5.2 cũng như tất cả các bước trong Mục 5.4, chúng tôi luôn bỏ qua các chuỗi tiền tố $\widetilde{0^{|f(x)|}} 0^{|f(x)|+1} 10^{|x|} 10^{p(|x|)} 1$ trong $\left|\phi^{p, f}(x)\right\rangle$ và $0^{|x|} 10^{p(|x|)} 1$ trong $\left|\phi^{p}(x)\right\rangle$, và chúng tôi chú ý đến các qubit còn lại. Hàm lượng tử mong muốn $g$ sẽ được xây dựng từng bước thông qua các Mục 5.2 5.5.

\subsection*{5.2 Xây dựng một cấu hình ban đầu đã mã hóa}
Cho một đầu vào nhị phân $x=x_{1} x_{2} \cdots x_{n}$, cấu hình ban đầu $\gamma_{0}$ của $M$ trên $x$ có dạng ( $q_{0}, 0, b \cdots b x b \cdots b$ ), và do đó mã $\left|\hat{\gamma}_{0}\right\rangle$ của $\gamma_{0}$ phải có dạng

$$
\left|q_{0}\right\rangle \otimes|\hat{3}, \hat{b}\rangle \otimes \cdots \otimes|\hat{3}, \hat{b}\rangle \otimes\left|\hat{2}, \hat{x}_{1}\right\rangle \otimes\left|\hat{3}, \hat{x}_{2}\right\rangle \otimes\left|\hat{3}, \hat{x}_{3}\right\rangle \otimes \cdots \otimes\left|\hat{3}, \hat{x}_{n}\right\rangle \otimes|\hat{3}, \hat{b}\rangle \otimes \cdots \otimes|\hat{3}, \hat{b}\rangle
$$

trong đó $\hat{x}_{i}$ là mã của $x_{i}, q_{0}\left(=0^{\ell}\right)$ là trạng thái nội ban đầu, và $x_{1}$ nằm ở ô 0 . Lưu ý rằng $\ell\left(\left|\hat{\gamma}_{0}\right\rangle\right)= 8 p(n)+\ell+4$. Trong phần tiếp theo, chúng tôi sẽ chỉ ra cách tạo mã cụ thể $\left|\hat{\gamma}_{0}\right\rangle$ này từ trạng thái lượng tử $\left|\phi^{p}(x)\right\rangle$. Để đơn giản, chúng tôi sẽ bỏ qua thuật ngữ $\left|\hat{q}_{0}\right\rangle$ trong các bước sau ngoại trừ Bước 8).

\begin{enumerate}
  \item Bắt đầu với đầu vào $\left|\phi^{p}(x)\right\rangle$, trước tiên chúng ta biến đổi nó thành $\left|0^{n} 1\right\rangle\left|0^{p(n)} 1\right\rangle\left|10^{11 p(n)+5} 1\right\rangle|x\rangle$ bằng một hàm lượng tử $h_{1}=\operatorname{Skip}[N O T]$, thỏa mãn cả hai $h_{1}\left(\left|0^{m} 1\right\rangle|\psi\rangle\right)=\left|0^{m} 1\right\rangle \otimes N O T(|\psi\rangle)$ và $h_{1}\left(\left|0^{m+1}\right\rangle\right)=\left|0^{m+1}\right\rangle$ cho bất kỳ số $m \in \mathbb{N}$ nào và bất kỳ trạng thái lượng tử $|\psi\rangle \in \mathcal{H}_{\infty}$. Bổ đề 3.9 đảm bảo rằng $h_{1}$ thực sự tồn tại trong $\widehat{\square_{1}^{\mathrm{QP}}}$.
\end{enumerate}

Từ $\left|0^{p(n)} 1\right\rangle\left|10^{11 p(n)+5} 1\right\rangle$, chúng ta muốn tạo ra $\left|0^{p(n)} 1\right\rangle\left|0^{p(n)} 1\right\rangle\left|0^{10 p(n)+5} 1\right\rangle$ bằng một hàm lượng tử $\widehat{\square_{1}^{\mathrm{QP}}}$ thích hợp $f_{1}$. Để định nghĩa hàm lượng tử mong muốn $f_{1}$, trước tiên chúng ta xây dựng một hàm lượng tử khác $g_{1}$ ánh xạ $\left|0^{p(n)} 1\right\rangle\left|0^{m} 10^{k} 1\right\rangle$ thành $\left|0^{p(n)} 1\right\rangle\left|0^{m+1} 10^{k-1} 1\right\rangle$ cho bất kỳ hai số nguyên $m \geq 0$ và $k \geq 1$ nào bằng cách đặt

$$
g_{1}(|\phi\rangle)= \begin{cases}|\phi\rangle & \text { if } \ell(|\phi\rangle) \leq 1 \\ |0\rangle \otimes g_{1}(\langle 0 \mid \phi\rangle)+|1\rangle \otimes g_{2}(\langle 1 \mid \phi\rangle) & \text { otherwise }\end{cases}
$$

trong đó $g_{2}$ được giới thiệu như

$$
g_{2}(|\phi\rangle)= \begin{cases}|\phi\rangle & \text { if } \ell(|\phi\rangle) \leq 1 \\ S W A P\left(|0\rangle \otimes g_{2}(\langle 0 \mid \phi\rangle)+|1\rangle\langle 1 \mid \phi\rangle\right) & \text { otherwise }\end{cases}
$$

Một cách chính thức hơn, ta đặt $g_{2}=Q R e c_{1}\left[I, S W A P, I \mid g_{2}, I\right]$ và đặt $\hat{g}_{2}=\operatorname{Branch}\left[I, g_{2}\right]$. Sau đó ta đặt $g_{1}= Q \operatorname{Rec}_{1}\left[I, I, \hat{g}_{2} \mid g_{1}, I\right]$. Hàm lượng tử $f_{1}$ cuối cùng được định nghĩa là $f_{1}=Q \operatorname{Rec}_{1}\left[I, g_{1}, I \mid f_{1}, I\right]$; cụ thể,

$$
f_{1}(|\phi\rangle)= \begin{cases}|\phi\rangle & \text { if } \ell(|\phi\rangle) \leq 1 \\ g_{1}\left(|0\rangle \otimes f_{1}(\langle 0 \mid \phi\rangle)+|1\rangle\langle 1 \mid \phi\rangle\right) & \text { otherwise }\end{cases}
$$

Theo cách tương tự, ta tiếp tục biến đổi $\left|0^{p(n)} 1\right\rangle\left|0^{10 p(n)+5} 1\right\rangle$ thành $\left|0^{p(n)} 1\right\rangle\left|0^{2 p(n)+2} 1\right\rangle\left|0^{8 p(n)+2} 1\right\rangle$. Sau đó ta đổi $\left|0^{n} 1\right\rangle\left|0^{2 p(n)+2} 1\right\rangle$ thành $\left|0^{n} 1\right\rangle\left|0^{n} 1\right\rangle\left|0^{2 p(n)-n+1} 1\right\rangle$ và $\left|0^{n} 1\right\rangle\left|0^{8 p(n)+2} 1\right\rangle$ thành $\left|0^{n} 1\right\rangle\left|0^{n} 1\right\rangle\left|0^{8 p(n)-n+1} 1\right\rangle$. Tổng thể, đầu vào $\left|\phi^{p}(x)\right\rangle$ được biến đổi thành $\left|0^{n} 1\right\rangle\left|\left(0^{p(n)} 1\right)^{3}\right\rangle\left|\left(0^{n} 1\right)^{2}\right\rangle\left|0^{2 p(n)-n+1} 1\right\rangle\left|0^{8 p(n)-n+1} 1\right\rangle|x\rangle$.\\
(*) Trong phần tiếp theo, bằng cách bỏ qua $\left|0^{n} 1\right\rangle\left|\left(0^{p(n)} 1\right)^{3}\right\rangle\left|\left(0^{n} 1\right)^{2}\right\rangle$, chúng ta tạm thời giả sử rằng đầu vào của chúng ta là $\left|0^{8 p(n)-n+1} 1\right\rangle|x\rangle$.\\
2) Để dễ đọc, chúng ta giải thích bước này bằng một ví dụ minh họa của $\left|0^{6} 1\right\rangle\left|x_{1} x_{2} x_{3}\right\rangle$, mà chúng ta muốn biến đổi thành $|00\rangle\left|\hat{3} \hat{x}_{1} \hat{x}_{2} \hat{x}_{3}\right\rangle$. Vì mục đích này, chúng ta bắt đầu bằng cách thay đổi $\left|0^{6} 1 x_{1} x_{2} x_{3}\right\rangle$ thành $\left|x_{3} x_{2} x_{1} 10^{6}\right\rangle$ bằng cách áp dụng REVERSE.

Để có được $\left|x_{3} 0 x_{2} 0 x_{1} 001\right\rangle|00\rangle$ từ $\left|x_{3} x_{2} x_{1}\right\rangle\left|10^{6}\right\rangle$, chúng ta tiếp tục áp dụng $g_{3}=S W A P \circ R E P_{1}$, biến đổi $\left|x_{3} x_{2} x_{1} 10^{6}\right\rangle$ thành $\left|x_{3} 0 x_{2} x_{1} 10^{5}\right\rangle$. Chúng ta lặp lại việc áp dụng $g_{3}$ và biến đổi $\left|x_{3} x_{2} x_{1}\right\rangle\left|10^{6}\right\rangle$ thành $\left|x_{3} 0 x_{2} 0 x_{1} 0\right\rangle\left|10^{3}\right\rangle$. Quá trình này có thể được thực hiện bằng hàm lượng tử $h_{3}=2 Q R e c_{2}\left[I, h^{\prime}, g_{3} \mid\left\{h_{s}^{\prime \prime}\right\}_{s \in\{0,1\}^{2}}\right]$ được định nghĩa bởi phép đệ quy lượng tử 2 qubit, trong đó $h^{\prime}=$ Branch $_{2}\left[\left\{h_{s}^{\prime}\right\}_{s \in\{0,1\}^{2}}\right]$ với $h_{a 1}^{\prime}=S W A P$ và $h_{a 0}^{\prime}=I$ cũng như $h_{a 1}^{\prime \prime}=I$ và $h_{a 0}^{\prime \prime}=h_{3}$ cho mỗi bit $a \in\{0,1\}$; cụ thể,

$$
h_{3}(|\phi\rangle)= \begin{cases}|\phi\rangle & \text { if } \ell(|\phi\rangle) \leq 2, \\ \sum_{a \in\{0,1\}}\left(S W A P\left(|a 1\rangle\left\langle a 1 \mid \psi_{g_{3}, \phi}\right\rangle\right)+|a 0\rangle \otimes h_{3}\left(\left\langle a 0 \mid \psi_{g_{3}, \phi}\right\rangle\right)\right) & \text { otherwise },\end{cases}
$$

trong đó $\left|\psi_{g_{3}, \phi}\right\rangle$ đại diện cho $g_{3}(|\phi\rangle)$. Chúng ta tiếp tục áp dụng $g_{4}=R E V E R S E \circ h_{3}$ để biến đổi $\left|x_{3} x_{2} x_{1}\right\rangle\left|10^{6}\right\rangle$ thành $\left|0^{3} 1\right\rangle\left|\hat{x}_{1} \hat{x}_{2} \hat{x}_{3}\right\rangle$. Trong trường hợp tổng quát, quá trình trên biến đổi $\left|0^{n} 1\right\rangle\left|x_{1} x_{2} \cdots x_{m}\right\rangle$ thành $|1\rangle\left|\hat{x}_{1} \hat{x}_{2} \cdots \hat{x}_{n}\right\rangle$.

Cuối cùng, chúng ta áp dụng $g_{5}$, biến đổi $\left|0^{3} 1\right\rangle\left|\hat{x}_{1} \hat{x}_{2} \hat{x}_{3}\right\rangle$ thành $|00\rangle\left|\hat{3} \hat{x}_{1} \hat{x}_{2} \hat{x}_{3}\right\rangle$, được định nghĩa bởi

$$
g_{5}(|\phi\rangle)= \begin{cases}|\phi\rangle & \text { if } \ell(|\phi\rangle) \leq 1 \\ S W A P\left(|00\rangle \otimes g_{5}(\langle 00 \mid \phi\rangle)+\sum_{y \in\{0,1\}^{2}-\{00\}}(|y\rangle\langle y \mid \phi\rangle)\right) & \text { otherwise }\end{cases}
$$


\end{document}